\documentclass[10pt,a4paper]{amsart}
\usepackage[margin=19mm]{geometry}
\usepackage{amsmath,amssymb,mathtools}
\usepackage[scr=rsfs,cal=euler]{mathalfa}
\usepackage{xfrac}
\usepackage[unicode,hidelinks,bookmarks=false]{hyperref}
\newcommand{\ie}{i.e.\ }
\newcommand{\cf}{cf.\ }
\newcommand{\resp}{respectively\ }
\newcommand{\Subsection}{\subsection}
\providecommand{\nobreakdash}{\nobreak\hbox{-}}
\setlength{\emergencystretch}{2em}
\multlinegap0pt
\usepackage{bm}
\usepackage{bbm}
\usepackage{dsfont}
\usepackage{xspace}
\usepackage{cancel}
\usepackage{mathtools}
\usepackage{amsthm}
\makeatletter
\@ifundefined{theorem}{\newtheorem{theorem}{Theorem}}{}
\@ifundefined{lemma}{\newtheorem{lemma}[theorem]{Lemma}}{}
\makeatother
\title[Rate estimates for total variation distance with application]{Rate estimates for total variation distance with applications}
\author{Miklos Rasonyi}
\date{Source: arXiv:2412.04013v2 (CC BY 4.0)}
\begin{document}
\maketitle

\noindent\emph{Source and licence.} Rate estimates for total variation distance with applications. Version arXiv:2412.04013v2 (CC BY 4.0), distributed under the Creative Commons Attribution 4.0 International licence, as recorded in the arXiv licence metadata for this exact version and shown on the abstract page https://arxiv.org/abs/2412.04013v2. Full rights notice: licenses/arxiv-2412.04013-CC-BY-4.0.txt.\par\smallskip

\noindent\textit{Editorial scope.} Counted unit: the Lemma whose source label is \texttt{bundi} in the article (source0.tex lines 840--845, source label \texttt{bundi}), delivered together with the article's own complete original proof of it (source0.tex lines 846--911, one \texttt{proof} environment of 66 lines). No mathematics is written, repaired or re-derived: statement and proof are the article's own text line for line. Required context retained uncounted: smoothie (lines 426--428); momenti (lines 430--432); sublemma (lines 913--915). Every definition, display and label of those lines is retained unchanged. Omitted from this package and not redistributed: 1--425, 429--429, 433--839, 912--912, 916--1176. Nothing inside a retained excerpt is omitted or altered: every retained body is delivered exactly as the article's lines are written, so no omission receipt of any kind accompanies this package. Citations: the retained excerpts carry no citation command at all, so no bibliography entry of the article is transcribed and no reference is redistributed. Reference adaptation: none was needed and none was made; every reference inside the delivered unit resolves inside it, so no unresolved reference or citation is produced. Printed slips kept literal: no printed word, symbol, hypothesis or number of the counted statement or of its proof was altered. Layout and class: the article's own class is \texttt{article} with packages including amsmath, amsthm, amssymb, latexsym, libertinus; the wrapper is the lane's pinned \texttt{amsart} editorial wrapper at 10pt on A4, which restates the theorem-like environments the retained text needs and adds the article's own macro definitions that the retained text needs (none), with extra wrapper packages (bm, bbm, dsfont, xspace, cancel, mathtools). The delivered numbering is the unit's own. Assets: no retained span contains a figure environment, a graphics-inclusion command, a PDF-page inclusion, a table or a TikZ picture; no image file is copied into this package, the package declares no asset dependency and no image file is redistributed with it. Licence: version arXiv:2412.04013v2 of this article is distributed under the Creative Commons Attribution 4.0 International licence, as recorded in the arXiv licence metadata for this exact version and shown on the abstract page https://arxiv.org/abs/2412.04013v2. A full rights notice is delivered at licenses/arxiv-2412.04013-CC-BY-4.0.txt. Characters: every retained excerpt is pure ASCII, so the delivered engine, XeLaTeX with its default Latin Modern text font, reports no missing character.
\begin{equation}\label{smoothie}
|\phi(u)|\leq C_{\sharp}(1+|u|)^{-\alpha},\ u\in\mathbb{R}^{d}
\end{equation}

\begin{equation}\label{momenti}
E[|Y_{0}|^{k}]<\infty.	
\end{equation}

\begin{lemma}\label{sublemma} Let \eqref{smoothie} hold. Then for each $c>0$,
there is $\rho=\rho(c)<1$ satisfying $\sup_{|u|\geq c}|\phi(u)|\leq \rho$.	
\end{lemma}

\begin{lemma}\label{bundi} If \eqref{smoothie} holds then
for each integer $l\geq 1$, there are $N(l),H_{l}>0$ such that
\begin{equation}\label{valami}
|\phi_{S_{n}}(u)|\leq H_{l}(1+|u|)^{-l},\ u\in\mathbb{R}^{d},\ n\geq N(l).
\end{equation}
\end{lemma}

\begin{proof}
Notice that $\phi_{S_{n}}(u)=\phi^{n}(u/\sqrt{n})$. We will present three separate arguments
depending on how large $u$ is. Fix an integer $l\geq 1$.

\smallskip

\noindent\textsl{Case 1: $|u|$ small.} As $Y_{0}$ has zero mean and unit covariance matrix and $\phi$ is smooth
by \eqref{momenti}, we have $\phi(0)=1$, $\phi'(0)=0$, $\phi''(0)=-\mathrm{Id}_{\mathbb{R}^{d}}$ 
and there is $0<c_{1}\leq 1$ such that 
$\phi(x)=1-|x|^{2}/2+O(|x|^{3})$ for $|x|\leq c_{1}$. Choosing an even smaller $c_{1}$ we may thus
guarantee that $|\phi(x)|\leq 1-|x|^{2}/3$ for all $|x|\leq c_{1}$. Recalling the Taylor expansion
of the logarithm, we may estimate 
$$
\left(1-\frac{|x|^{2}}{3}\right)\leq e^{-|x|^{2}/4}
$$
for all $|x|\leq c_{2}$, with $c_{2}\leq c_{1}$ small enough, hence
$$
|\phi_{S_{n}}(u)|\leq \left(1-\frac{|u|^{2}}{3n}\right)^{n}\leq e^{-|u|^{2}/4}
$$
for all $n$ and for all $|u|\leq c_{2}\sqrt{n}$.

Now notice that $(1+x)^{l}\leq l!e^{x+1}\leq l! e^{\frac{x^{2}}{4} +3}$ for all $x\geq 0$,
hence also 
$$
|\phi_{S_{n}}(u)|\leq \frac{e^{3}l!}{(1+|u|)^{l}},
$$
showing \eqref{valami} for $|u|\leq c_{2}\sqrt{n}$.

\smallskip{}

\noindent\textsl{Case 2: $|u|$ large.} We may and will assume $C_{\sharp}\geq 1$. By \eqref{smoothie}, one has 
\begin{equation}\label{tavol}
|\phi(x)|\leq C_{\sharp}|x|^{-\alpha},\ x\in\mathbb{R}^{d}.	
\end{equation} 
Let $n\geq \hat{N}(l):=2l/\alpha$. 
Set $J:=C_{\sharp}^{2/\alpha}$. Notice that for $|u|\geq Jn$, one has
$$
|u|\geq n^{\frac{\alpha n}{2(\alpha n -l)}}C_{\sharp}^{\frac{n}{\alpha n-l}}
$$
but then, since \eqref{tavol} holds,
$$
|\phi_{S_{n}}(u)|\leq \frac{(\sqrt{n})^{\alpha n}C_{\sharp}^{n}}{|u|^{\alpha n} }\leq \frac{1}{|u|^{l}}.
$$
This implies 
$$
|\phi_{S_{n}}(u)|\leq \frac{2^{l}}{(1+|u|)^{l}}
$$
since $|u|\geq Jn\geq 1$.

\smallskip{}


\noindent\textsl{Case 3: $|u|$ moderate.}
Apply Lemma \ref{sublemma} below
with the choice $c:=c_{2}$. We obtain $\rho$ such that $|\phi_{S_{n}}(u)|\leq \rho^{n}$
for all $|u|\geq c_{2}\sqrt{n}$. 
It is clear that $\rho^{n}\leq (1+Jn)^{-l}$ holds for $n\geq \tilde{N}$ with some 
$\tilde{N}=\tilde{N}(l)$.
This shows \eqref{valami} for all $c_{2}\sqrt{n}\leq |u|\leq Jn$.

\smallskip


We have thus verified \eqref{valami} for all $l\geq 1$, $n\geq N(l):=\max\{\tilde{N}(l),\hat{N}(l)\}$ 
and $u\in\mathbb{R}^{d}$ and may conclude. 
\end{proof} 

\end{document}
